我们继续关于范畴迹的一般讨论。设 $\mC$ 如命题 \ref{L:trace cat} 中那样。
先暂且假设 $\sigma$ 是恒等函子。此时把 $\Tr_\sigma(\mC)$ 记作 $\Tr(\mC)$。对 $\mC$ 中每个自同态 $f: V\to V$，记 $\coev_f:\mathbf{1}\to V\otimes V^*$ 为由 $\mathbf{1}\xrightarrow{\coev_V}V\otimes V^*\xrightarrow{f\otimes\id_{V^*}}V\otimes V^*$ 给出的态射。则有元素
\[S_{\coev_f,\ev_V}\in \End_{\Tr(\mC)}\widetilde{\mathbf{1}}.\]
特别地，当 $f=\id$ 时，记
\begin{equation}
\label{E:abs S}
S_V:=S_{\coev_{V},\ev_{V}}\in \End_{\Tr(\mC)}\widetilde{\mathbf{1}}.
\end{equation}
该元素只依赖于 $V$ 的同构类。由构造可得
\[S_{\mathbf{1}}=\id,\quad S_{V'}\cdot S_{V}=S_{V\otimes V'},\quad S_{V\oplus V'}=S_V+S_{V'}.\]
于是若 $\mC$ 是加性的，就有典范同态
\begin{equation}
\label{E: K to End}
K^{\oplus}(\mC)^{\on{op}}\to \End_{\Tr(\mC)}\widetilde{\mathbf{1}},
\end{equation}
其中 $K^{\oplus}(\mC)$ 表示 $\mC$ 的分裂 Grothendieck 环：作为 Abel 群，它由 $\mC$ 中对象生成，关系为 $[X]+[Y]=[X\oplus Y]$；环结构由张量积给出。

\begin{rmk}
注意若 $\mC$ 是对称幺半的，则 $\mC$ 的恒等函子由交换性约束给出自然的迹函子结构，从而分解为 $\mC\to \Tr(\mC)\to \mC$。显然诱导映射 $\End_{\Tr(\mC)}(\widetilde{\mathbf{1}})\to \End_\mC(\mathbf{1})$ 把 $S_{\coev_f,\ev_V}$ 送到 $f$ 的通常迹 $\tr_V(f)$，把 $S_V$ 送到 $V$ 的秩 $\on{rank}(V)$（这些概念的定义见例如 \cite[(1.7.3)]{DM}）。
\end{rmk}

上述构造有如下推广。设 $X\in \mC$，$V\in \mC$，并给定同构 $c_{X,V}: X\otimes V\simeq V\otimes X$。则由
\[X\xrightarrow{\coev_V\otimes\id_X}V\otimes V^*\otimes X, \quad V^*\otimes X\otimes V\xrightarrow{\id_{V^*}\otimes c_{X,V}} V^*\otimes V\otimes X\xrightarrow{\ev_V\otimes \id_X} X\]
得到元素 $S_{c_{X,V}}\in \End_{\Tr(\mC)}\widetilde{X}$。
若把 $X$ 在 $\mC$ 中的中心化范畴定义为这样的范畴 $\mZ_{\mC}(X)$：对象是偶 $\{V\in\mC, c_{X,V}: X\otimes V\simeq V\otimes X\}$，态射是 $\mC$ 中与 $c_{X,V}$ 和 $c_{X,V'}$ 相容的态射 $V\to V'$；则 $\mZ_{\mC}(X)$ 是幺半范畴，且当 $\mC$ 还加性时 $\mZ_{\mC}(X)$ 也加性。于是有态射
\begin{equation}
\label{E: K to End2}
K^{\oplus}(\mZ_{\mC}(X))^{\on{op}}\to \End_{\Tr(\mC)}(\widetilde X).
\end{equation}
特别地，若 $\mC$ 是加性双幺半范畴，则对每个 $X$ 都有自然幺半函子 $\mC\to \mZ_\mC(X)$。因此此时有
$$K^{\oplus}(\mC)^{\on{op}}\to \End_{\Tr(\mC)}\widetilde{X}.$$

现在处理一般的 $\sigma$。记 $\mC^\sigma$ 为 $\mC$ 中 $\sigma$-等变对象构成的范畴，即对象是偶 $(X,\phi)$，其中 $X\in\mC$，$\phi: \sigma X\simeq X$ 是同构；$(X,\phi)$ 与 $(X',\phi')$ 之间的态射是 $\mC$ 中与 $\phi$ 和 $\phi'$ 相容的态射。注意 $\mC^\sigma$ 自然地是幺半范畴，且若 $\mC$ 中每个对象都有左对偶，则 $\mC^\sigma$ 中亦然。
于是可以把上述讨论用于 $\mC^\sigma$。遗忘函子 $\mC^\sigma\to\mC$ 诱导函子
\[\Tr(\mC^\sigma)\to \Tr_\sigma(\mC).\]
特别地有
\begin{equation}
\label{E: K to End3}
K^{\oplus}(\mC^\sigma)^{\on{op}}\to \End_{\Tr_\sigma(\mC)}(\widetilde{\mathbf 1}).
\end{equation}
具体地，对 $(X,\phi: \sigma X\simeq X)\in\mC^\sigma$，$\widetilde{\mathbf 1}$ 在 $\Tr_\sigma(\mC)$ 中对应的自同态由 $S_{\coev_X, \ev_X(1\otimes\phi)}$ 给出。

下面假设 $E$ 是代数闭域，$\mC$ 是半单 Abel 范畴，且单位对象 $\mathbf{1}$ 不可约。则 $\mC^\sigma$ 是 Abel 幺半范畴。

\begin{prop}
\label{L: K iso End}
在上述假设下，映射 \eqref{E: K to End3} 诱导同构
\[
(K(\mC^\sigma)^{\on{op}}\otimes E)/I\cong \End_{\Tr_\sigma(\mC)}(\widetilde{\mathbf 1}),
\]
其中 $I$ 是 $K(\mC^\sigma)^{\on{op}}\otimes E$ 中由 $[(X,a\phi)]-a[(X,\phi)]$ 生成的理想，$(X,\phi)$ 取遍 $\mC^\sigma$ 中的不可约对象，$a\in E^\times$。
\end{prop}
结合下面的注 \ref{R: cat center}，Ostrik 有类似的结果 \cite[推论 2.16]{O}。
\begin{proof}
由于该结果后文不用，我们只给出证明梗概。首先，因 $E$ 代数闭且 $\mathbf{1}$ 不可约，可以说明：若对某个 $u:\mathbf{1}\to V\otimes W$ 与 $v: W\otimes \sigma V\to \mathbf{1}$ 有 $S_{u,v}\neq 0$，则 $S_{u,v}$ 是 $S_{\coev_{Y}, \ev_Y(1\otimes\psi)}$ 的线性组合，其中 $(Y,\psi)$ 不可分解。
其次可以说明，$\mC^\sigma$ 中的不可分解对象都形如 $(Y,\psi)=(X\otimes W, \phi\otimes f)$，其中 $(X,\phi)$ 在 $\mC^\sigma$ 中不可约，$W$ 是有限维 $E$-向量空间，$f$ 是 $W$ 的自同构。记 $f^{\on{ss}}$ 为 $f$ 的半单化，$(Y,\psi^{\on{ss}})=(X\otimes W,\phi\otimes f^{\on{ss}})$ 为 $(Y,\psi)$ 的半单化。
再说明 $S_{\coev_{Y}, \ev_Y(1\otimes\psi)}=S_{\coev_{Y},\ev_Y(1\otimes\psi^{\on{ss}})}$。由此 \eqref{E: K to End3} 诱导出所要的映射，且它是满射。单射性显然。
\end{proof}

\begin{rmk}回忆传统的范畴化/去范畴化是借助 $K$-环构造在幺半范畴与代数之间建立通道。上述命题提示了另一条经由范畴迹构造的通道。
\end{rmk}

\begin{rmk}
\label{R: cat center}
幺半范畴还有范畴中心（或称 Drinfeld 中心）的概念（见 \cite{BFO, BN, Lu3}）。给定 $(\mC,\otimes,\mathbf{1})$，其范畴中心 $\mZ(\mC)$ 是这样一个范畴：对象由 $(X\in \mC, \{\al_{X,Y}: X\otimes Y\simeq Y\otimes X\}_{Y\in \mC})$ 组成，并满足自然的相容条件。在某些情形下，有等价 $\mZ(\mC)\to \Tr(\mC)$。不过在我们的应用中，使用范畴迹的概念更为自然。
\end{rmk}
